% Generated draft: requires mathematical and source review.
\documentclass[11pt]{article}
\usepackage[margin=25mm]{geometry}
\usepackage{fontspec}
\setmainfont{Latin Modern Roman}
\usepackage{amsmath,amssymb,mathtools,unicode-math}
\setmathfont{Latin Modern Math}
\usepackage{enumitem,needspace,xurl,hyperref}
\setlength{\parindent}{0pt}
\setlength{\parskip}{5pt}
\setlength{\emergencystretch}{3em}
\newcommand{\N}{\mathbb{N}}
\newcommand{\Z}{\mathbb{Z}}
\newcommand{\Q}{\mathbb{Q}}
\newcommand{\R}{\mathbb{R}}
\newcommand{\C}{\mathbb{C}}
\newcommand{\F}{\mathbb{F}}
\newcommand{\A}{\mathbb{A}}
\newcommand{\PP}{\mathbb{P}}
\newcommand{\E}{\mathbb{E}}
\newcommand{\eps}{\varepsilon}
\newcommand{\dd}{\,\mathrm{d}}
\begin{document}
\section*{UnsolvedMath ID: 20002839}
\textbf{Source problem code:} catalogue entry
\subsection{Definitions and mathematical statement}
Let $\bar\rho$ be the residual two-dimensional Galois representation implicit in the workshop's deformation problem, let $R=R^{\mathrm{mod}}$ be its indicated modular deformation ring, let $I\subset R$ be the indicated ideal, put $A=R/I$, and let $\rho_A$ be the universal deformation over $A$. These rings, the ideal, tame level and completed-cohomology localization are referenced but not fully defined in this excerpt. For the finite set of places $\Sigma$ containing $p$, Conjecture 4 displays an isomorphism of the form
\[
\widehat H^1[I]\ \stackrel{?}{\simeq}\
\rho_A\widehat\otimes_A
\widehat\pi_p(\rho_A|_{G_{\Q_p}})
\widehat\otimes_A
\left(\widehat\bigotimes_{\substack{\ell\in\Sigma\\\ell\ne p},A}
\widehat\pi^{\mathrm{mod}}_\ell(\rho_A|_{G_{\Q_\ell}})\right),
\]
equivariant for $G_\Q\times\prod_{q\in\Sigma}\operatorname{GL}_2(\Q_q)$. The Galois factor has the $G_\Q$ action, and the local factors have the corresponding general linear group actions. Away from $p$, the source describes the factor as the $\mathfrak m$-adic completion of modified local Langlands on the generic fiber, descended to $A$. The topology and descent construction need specification. The original left-hand notation contains ``??''; its meaning is not silently corrected.

The source's Theorem 14 says that the display holds under its two hypotheses. First, the residual local representation has scalar endomorphisms; the compressed wording indicates the irreducible case or a nonsplit reducible case
\[
\bar\rho|_{G_{\Q_p}}\simeq
\begin{pmatrix}\chi&*\\0&\psi\end{pmatrix},\qquad \chi\ne\psi,\quad *\ne0,
\]
excluding a twist of $\begin{pmatrix}\omega^{-1}&*\\0&1\end{pmatrix}$, with $\omega$ the mod-$p$ cyclotomic character. The exact alternative phrasing requires source review. Second, there exists the stated orthonormalizable, $\mathfrak m$-adically complete $A$-module $\widehat\pi_p(A)$ interpolating the classical Banach representations: for every classical modular form $f$ of weight $k\geq2$ deforming $\bar\rho$, unramified outside $\Sigma\setminus\{p\}$ in the source's wording, whose map $\varphi_f:R\to K_f$ factors through $A$,
\[
\widehat\pi_p(A)\widehat\otimes_{A,\varphi_f}K_f
\simeq B(\rho_f|_{G_{\Q_p}}).
\]
The source attaches Berger--Breuil--Emerton to this interpolation. Whether the ramification wording refers to tame ramification and the exact Banach/integral conventions remain unresolved.

The accompanying Remark 1 says that then-unproved Colmez results could remove the second hypothesis. Corollary 15, under the same hypotheses and for a deformation $\rho:G_\Q\to\operatorname{GL}_2(E)$ with $R\to E$ factoring through $A$, states: if $\rho|_{G_{\Q_p}}$ is potentially semistable and trianguline with distinct Hodge--Tate weights, then $\rho$ arises from a classical modular form; if it is trianguline, then $\rho$ comes from a twist of an overconvergent finite-slope eigenform with tame level equal to the tame conductor of $\rho$. Remark 2 reports that the trianguline condition in the first assertion should be removable according to correspondence with Colmez, relates this to the $\operatorname{GL}_2$ Fontaine--Mazur conjecture and Kisin's talk, and claims a two-variable $p$-adic $L$-function over the part of the eigencurve mapping to $A$, as well as an unspecified mod-$p$ statement. These are historical assertions with the source's qualifications.

A separate algebraic issue identified by the supplied aid concerns the left side. If brackets mean the conventional annihilator submodule, define
\[
P[I]=\{x\in P:Ix=0\},\qquad C_I=\operatorname{Ann}_R(I).
\]
For every finite projective $R$-module $P$, the canonical multiplication map gives
\[
C_I\otimes_R P\simeq P[I],
\]
whereas quotient base change gives $P/IP=A\otimes_RP$. These generally differ. The aid's reduced complete local example, over any field $k$, is
\[
R=k[[x,y,z]]/(xy,xz),\qquad I=(x),\qquad
A=k[[y,z]],\qquad C_I=(y,z).
\]
As an $A$-module, $C_I$ requires two generators at the maximal ideal and is not an invertible rank-one module. Thus reducedness alone does not identify the annihilator with the quotient. Deriving a component formula from a full-ring factorization needs the congruence factor $C_I$, or hypotheses permitting it to be trivialized or absorbed into a multiplicity module. Extension to completed, non-finite-projective factors additionally requires the relevant flatness, saturation, inverse-limit and topology conditions. The aid mentions an adically compatible finite-level criterion and a dense-fiber counterexample, but does not supply their statements or hypotheses; neither is inferred from pointwise classical interpolation alone.

The second extracted remark continues into Section 11 (Schein's coefficient systems and diagrams) and Section 12 (Taylor on Herzig's thesis). These are later-topic extraction material, not extra assumptions of Conjecture 4.
\subsection{Short English statement}
Formulate the workshop's conditional completed-cohomology factorization and its modularity consequences, while retaining the uncertainty in $[I]$. If $[I]$ means annihilation by $I$, component extraction introduces $\operatorname{Ann}_R(I)$ and need not agree with quotient base change.
\subsection{Sources}
\begin{itemize}
\item Local UnsolvedMath registry, permanent ID $20002839$, source code \texttt{catalogue entry}; canonical \texttt{aim-representation-theory-notes.json}, \texttt{notes[110]}.
\item Michael Volpato's notes, AIM workshop \emph{$p$-adic representations, modularity, and beyond}, February 20--24, 2006: Conjecture 4, display (3), Theorem 14, Remarks 1--2 and Corollary 15; \url{https://aimath.org/WWN/padicmodularity/padicmodularity.pdf}. The supplied verification says the notes warn that they were typeset during talks and may contain transcription errors. No original literature field is supplied.
\item The source names Berger--Breuil--Emerton, Colmez and Kisin as above. The research aid also refers to later refined full-Hecke factorizations in broad generic settings, without an authoritative complete citation or exact hypotheses.
\item Research attempt 1: \texttt{partial\_result}, status \texttt{partially\_solved}, classification \texttt{PARTIAL-PROGRESS}, rigor \texttt{proved}, novelty confidence low. Its congruence-module obstruction and further criteria are machine-generated aid, not peer reviewed, and require independent verification before citation.
\end{itemize}
\end{document}
